\chapter*{Resumen}

\section*{\tituloPortadaVal}
Este trabajo realiza una comparativa de distintos modelos y técnicas de IA que logren replicar y hacer frente a la simulación de físicas costosas en proyectos 3D en motores de videojuegos. \newline

Se toma como ejemplo el caso de las simulaciones de telas, cuya geometría y caracteríticas, y consecuente alto coste computacional en simulación, lo convierte en campo experimental para la optimización de físicas mediante redes de neuronas. Para lograr simular una tela, se crea una simulación en Unity3D con un objeto con componente Cloth y un colisionador que genera movimieto y deformaciones en la malla. En el mismo entorno de Unity se generan los datasets, los que se exportan al entorno de entrenamiento para analizar diferentes modelos en Pytorch: una MLP, una RNN, y ténicas de propagación de error así como reducciones de espacio.

Las redes entrenadas se exportan de vuelta a Unity, y en este se ponen en marcha los modelos para simular el movimiento (mediante el paquete Sentis). Se modifican los vértices de una tela con un componente propio ML, que pide al modelo los deesplazamientos a realizar para cada uno. Finalmente, se realizan experimentos y comparativas entre los distintos modelos, tanto en cuanto a precisión del modelo como a resultados visuales. Se concluye que el modelo que obtiene mejores resultados replicando visualmente la tela fue una red neuronal recurrente con una técnica de propagación de error a futuro, no obstante el rendimiento aún precisa de optimizaciones para ser prometedor.\newline



\section*{Palabras clave}

\noindent Inteligencia Artificial, Simulación física, Unity, Videojuegos, Optimización, PyTorch, ML




